% Source: reports/experiments/main/summary.json (실험 current: 원본과 완성본을 ODT로 전처리, docs/experiment.md)
% created_at: 2026-10-18T22:09:46Z; run_id: main; unit: seconds; tasks: 321
% GPT-6-Luna = png_150__gpt6luna_odt, Docling = vector_pdf__docling_odt (ODT 직접 출력); MinerU2.5는 논문 비교 대상에서 제외.
% Columns: labels.*.{ssim_blur_native_pages_median,ratio_median,extra_pages_share}
% 시각 유사도·편집 시간(ratio, 원본 대비 배수)은 개별 평가값의 중앙값(소수 두 자리), 페이지 넘침은 출력본보다 페이지 수가 늘어난 복원본의 비율(백분율 한 자리).
% 편집 시간은 원본 321개에 대해 실패를 포함한 결과; 시각 유사도와 페이지 넘침은 시각 평가가 있는 성공 건 기준.
% 양성 대조군 먼저, 실제 방식은 ratio_median 오름차순, 음성 대조군 마지막.
\begin{table}[!htb]
\centering
\caption{복원 방식별 평가 결과}
\label{tab:main-results}
\papertablestyle
\begin{tabularx}{\linewidth}{@{}Lrrr@{}}
\toprule
방식 & \begin{tabular}[b]{@{}r@{}}시각 유사도\\(최대 1)\end{tabular} & \begin{tabular}[b]{@{}r@{}}편집 시간\\(배)\end{tabular} & \begin{tabular}[b]{@{}r@{}}페이지 넘침\\(\%)\end{tabular} \\
\midrule
원본 & 1.00 & 1.00 & 0.0 \\
\midrule
Word & 0.73 & \textbf{2.34} & 26.8 \\
한/글 & 0.71 & 2.49 & 17.1 \\
GPT-6-Luna & 0.73 & 2.67 & 34.9 \\
Docling & 0.71 & 2.77 & 20.9 \\
\midrule
이미지 & 0.99 & 3.51 & 0.0 \\
\bottomrule
\end{tabularx}
\end{table}
